\chapter{Fazit}

Das Einstellen der Hochspannung und die Beobachtung der Signale am Szintillator funkionierte gut und zeigte die erwartete Signalform. Die Beobachtung der Signale des HPGe Detektors lieferte ebenfalls gut erkennbare Signale in der erwarteten Form.

Anschließend wurden Spektren von Cobalt-60, Cäsium137 und Europium-154 aufgenommen, an denen eine Energiekalibration für die beiden Detektoren, siehe Gl \eqref{eq:cal_ge} und \eqref{eq:cal_nai}, durchgeführt wurde. Die Kalibration des HPGe Detektors kann mit einer Anpassungsgüte von $0.999$ also sehr gut betrachtet werden und hat im folgenden gut funktioniert. Die Kalibration des NaI Detektors hat ebenso gut funktioniert war aber für die weitere Durchführung weniger relevant.

Bei der Betrachtung der Halbwertsbreiten wurde festgestellt, dass der NaI Szintillator stets eine größere Linienbreite als der HPGe Detektor hat und bei beiden die Linienbreite mit zunehmender Energie steigt. Dies entspricht genau den Erwartungen. Anschließend wurde die Linienbreite des HPGe Detektors untersucht und die Anteile der intrinsischen und elektronischen Halbwertsbreiten quantifiziert. \todo{maybe more??}

Es wurden die Detektoreffizienzen anhand von Cäsium für beide Detektoren bestimmt, wobei der Szintillator erwartungsgemäß mit $\epsilon_\text{NaI} = 0.101(12)$ ($\approx 10.1\%$) deutlich über der des HPGe Detektors mit $\epsilon_\text{HPGE} = \num{0.0567(68)}$ ($\approx 5.7\%$) liegt. 
Die Energieabhängigkeit der Effizienz wurde für den HPGe Detektor quantifiziert und diese Anpassung mit der Effizienz von Cäsium überprüft. Es war eine gute übereinstimmung gegeben (siehe Abb \ref{fig:ge_efficiency}).
Mittels der bestimmten Effizienzfunktion wurde die Aktivität der Cobalt Probe abgeschätzt, wobei eine starke Abweichung aufgetreten ist, welche vermutlich von dem sehr kleinen Abstand der Cobalt Probe zum Detektor kommt.


Zuletzt wurde eine Langzeitmessung einer Bodenprobe und eine Langzeitmessung des Untergrundes durchgeführt. Die meisten gefundenen Gammalinien konnten plausiblen Radioisotopen zugeordnet werden. Einzig eine Linie bei ca. $\SI{1.9}{\mega\electronvolt}$ konnte keinem sinnvollen Isotop zugeordnet werden. Es konnte natürlich vorkommendes Kalim-40 sowie Zivilisatorisch entstandenes Cäsium-137 gefunden werden. Weiterhin wurden kosmogene Nuklide (Natrium-22 und Beryllium-7) gemessen werden, sowie Produkte aus der Zerfallsreihe von Uran, wie Blei-214 und Bismuth-214. Diese Elemente sind in einer Bodenprobe aus der Umgebung zu erwarten.
